\section{Signed periodic responses and harmonic defects}
\label{app:periodic}

This appendix proves the periodic-response and harmonic-defect
estimates used in the global single-parity theorem.

For $a>0$, define the paired series
\begin{equation}\label{eq:periodic-response}
 T_a(b)=\tanh b+\sum_{j\ge1}(-1)^j
       [\tanh(b+2aj)+\tanh(b-2aj)].
\end{equation}
It and every fixed derivative converge locally uniformly for $a>0$.
The function is odd and satisfies $T_a(b+2a)=-T_a(b)$.

\begin{theorem}[The exact signed periodic maximum]
\label{thm:periodic-max}
For every $a>0$,
\[
 T_a'(b)>0\quad(-a<b<a),\qquad T_a'(a)=0,\qquad T_a''(a)<0.
\]
Consequently
\[
 |T_a(b)|\le T_a(a)=F_0(a),
\]
with equality exactly at $b\in(2\mathbb Z+1)a$. There is an absolute
$c>0$ such that, for $a\in[1,2]$ and every real $b$,
\begin{equation}\label{eq:periodic-quadratic}
 F_0(a)-|T_a(b)|\ge
 c\,\operatorname{dist}(b,(2\mathbb Z+1)a)^2.
\end{equation}
\end{theorem}
\begin{proof}
Euler's product for $\cosh$ gives
\[
 \sech^2x=\prod_{j\ge0}(1+x^2/\lambda_j)^{-2},\qquad
 \lambda_j=((2j+1)\pi/2)^2.
\]
Let $S_j$ be independent gamma variables of shape two and rate
$\lambda_j$. Their sum $S$ is finite and positive almost surely,
because $\sum_j\E S_j<\infty$. Their Laplace transforms give
\begin{equation}\label{eq:sech-gaussian-mixture}
 \sech^2x=\E e^{-Sx^2}.
\end{equation}
For fixed $s>0$, the alternating Gaussian sum
\[
 G_{a,s}(b)=\sum_{j\in\mathbb Z}(-1)^j e^{-s(b+2aj)^2}
\]
is, after multiplication by $\sqrt{s/\pi}$, the Dirichlet heat kernel
on $(-a,a)$, started at zero, at time $1/(4s)$. The method-of-images
formula gives this identity directly. Positivity of the heat kernel
and the boundary point lemma imply
\[
 G_{a,s}(b)>0\ (|b|<a),\qquad G_{a,s}(a)=0,\qquad
 G_{a,s}'(a)<0.
\]
On the other hand,
\[
 T_a'(b)=\sum_{j\in\mathbb Z}(-1)^j\sech^2(b+2aj)
         =\E G_{a,S}(b).
\]
Interchange of sum and expectation is justified by the finite sum of
absolute expectations $\sum_j\sech^2(b+2aj)$. At the boundary,
differentiation is also justified: the absolute expectations of the
differentiated terms sum to
$\sum_j| (\sech^2)'((2j+1)a)|<\infty$.
Thus $T_a''(a)=\E G_{a,S}'(a)<0$.

Oddness and anti-periodicity prove the maximum statement. Substituting
$b=a$ into \eqref{eq:periodic-response} and telescoping the exponentially
small tails of $\tanh x=1-2/(1+e^{2x})$ gives
\[
 T_a(a)=1-4\sum_{j\ge0}\frac{(-1)^j}{1+e^{2a(2j+1)}}=F_0(a).
\]
Finally, the deficit in \eqref{eq:periodic-quadratic} has a strictly
positive quadratic coefficient at each maximum. Away from the maxima
it is strictly positive. Compactness in $a\in[1,2]$ and in the
normalized periodic phase proves the uniform bound.
\end{proof}

\begin{theorem}[Stability of an arbitrary number of phase defects]
\label{thm:phase-stability}
Let $k\ge1$, let $\eta\in\{0,1\}$ with $k+\eta$ odd, and let
$b_1,\ldots,b_k\in\R$. For $a\in[1,2]$, put
\[
 H=\E_Y\left[\chi_{[k]}(Y)
 T_a\left(\eta a+\sum_i b_iY_i\right)\right].
\]
There is an absolute $c>0$, independent of $k$, such that
\begin{equation}\label{eq:phase-stability}
 F_0(a)-|H|\ge c\min\left\{
 \sum_i\operatorname{dist}(b_i,(2\mathbb Z+1)a)^2,1\right\}.
\end{equation}
The extremal equality $|H|=F_0(a)$ holds exactly when every $b_i$ is an odd integer multiple of $a$.
\end{theorem}
\begin{proof}
Choose odd integers $r_i$ with $d_i=b_i-r_i a\in[-a,a]$.
Anti-periodicity cancels the parity character, giving, up to a single
common sign,
\[
 H=\E T_a\left(a+\sum_i d_iY_i\right).
\]
Theorem~\ref{thm:periodic-max}, including its strict second derivative,
implies the uniform pointwise inequalities
\[
 F_0(a)\mp T_a(a+z)\ge c_0
       [1\mp\cos(\pi z/(2a))].
\]
Indeed the ratios extend positively across the appropriate maxima and
are positive and periodic elsewhere. Taking expectations, the cosine
average is
\[
 p=\prod_i\cos(\pi d_i/(2a))\in[0,1].
\]
Both signed deficits are therefore at least $c_0(1-p)$. Since
$\cos t\le e^{-t^2/2}$ on $[-\pi/2,\pi/2]$,
\[
 1-p\ge1-\exp\left(-\frac{\pi^2}{8a^2}\sum_i d_i^2\right)
 \ge c_1\min\left\{\sum_i d_i^2,1\right\}.
\]
This proves the claim and the equality classification.
\end{proof}

\begin{lemma}[Restoring curvature under phase averaging]
\label{lem:periodic-curvature}
For $a\in[1,2]$, write $F=F_0(a)$ and $U(z)=T_a(a+z)$.
There are positive parameters $A$ and $\mu<3/4$, depending on $a$,
such that
\begin{equation}\label{eq:restoring-curvature}
 U''(z)\le-A(1-\mu)F+A(1-\mu/2)(F-U(z)).
\end{equation}
In particular, for any real random variable $Z$,
\[
 \E U(Z)\ge\tfrac45F
 \quad\Longrightarrow\quad
 \E U''(Z)\le-\tfrac12 A(1-\mu)F<0.
\]
\end{lemma}
\begin{proof}
Let the Jacobi elliptic modulus $k$ have nome
$q=e^{-\pi^2/(2a)}$, and let $K$ be its complete elliptic integral.
The Fourier series for the Jacobi function $\operatorname{sn}$ gives
\[
 T_a(b)=F_0(a)\operatorname{sn}(Kb/a,k),\qquad F_0(a)=kK/a.
\]
This follows by comparison with
\[
 T_a(b)=\frac\pi a\sum_{j\ge0}
 \frac{\sin((2j+1)\pi b/(2a))}
 {\sinh((2j+1)\pi^2/(4a))}.
\]
Set $A=(K/a)^2$ and $\mu=k^2$. The differential equation for
$\operatorname{sn}$ is
\[
 U''=-A(1+\mu)U+2A\mu U^3/F^2.
\]
For $x=U/F\in[-1,1]$, subtracting the right side of
\eqref{eq:restoring-curvature} from this expression gives
$-AF\mu(1-x)(2x+1)^2/2\le0$.

For completeness, $q<17/200$ on the indicated gain interval. The theta
series imply
\[
 \mu=\left(\frac{\vartheta_2(0,q)}{\vartheta_3(0,q)}\right)^4
 \le\frac{16q[1+q^2/(1-q^4)]^4}{(1+2q)^4}<\frac34.
\]
The rational upper bound increases for $q\le17/200$ and is less than
$3/4$ at $17/200$. The bound on $q$ follows from
$\pi>3141/1000$ and a positive Taylor lower bound for the exponential.
Finally,
$(1-\mu)/(1-\mu/2)\ge2/5$, which proves the expectation assertion.
\end{proof}

\subsection{Exact harmonic vectors}

The following argument treats a growing number of odd harmonics,
including harmonics of unbounded size.

\begin{theorem}[Odd harmonics cannot improve the balanced maximum]
\label{thm:all-harmonics}
For all sufficiently large odd $n$, every $a\in[1,2]$, and every
vector of positive odd integers $r_1,\ldots,r_n$,
\[
 \frac{\left|\E\chi_{[n]}(X)
       \tanh(a\sum_i r_iX_i)\right|}
 {a\sqrt{\sum_i r_i^2}}
 \le\frac{q_n(a_n)}{a_n\sqrt n}.
\]
Equality holds exactly when every $r_i=1$ and $a=a_n$.
The same statement allows arbitrary signs of the coefficients.
\end{theorem}
\begin{proof}
Let $k$ count the $r_i$ different from one and put
$\Delta=\sum_i(r_i^2-1)\ge8k$. The case $k=0$ is the balanced-gain
 theorem. Suppose first $1\le k\le n/2$.
The Fourier representation
\[
 \tanh x=\int_0^\infty\frac{\sin(\omega x)}{\sinh(\pi\omega/2)},d\omega
\]
can be folded onto $(0,\pi)$ after setting $\theta=a\omega$.
Because every $r_i$ and $n$ are odd, the product of sines changes sign
between consecutive half-periods. The folded kernel is
\[
 K_a(\theta)=\sum_{j\ge0}\frac{(-1)^j}
 {\sinh(\pi(\theta+j\pi)/(2a))}>0.
\]
The alternating-series test proves its strict positivity. For real
$p>0$, define
\[
 J_p(a)=\int_0^\pi\sin^p\theta\,K_a(\theta)\,d\theta.
\]
For odd $n$, $q_n(a)=J_n(a)/a$. Dropping the absolute values of the
$k$ exceptional sine factors bounds the weighted coefficient by
$J_{n-k}(a)/a$.

Uniformly for $a\in[1,2]$, Laplace expansion at $\theta=\pi/2$ gives
\begin{equation}\label{eq:folded-laplace}
 J_p(a)=\sqrt{2\pi}\,K_a(\pi/2)p^{-1/2}
 \left[1+\frac{b_a}{p}+O(p^{-2})\right],\qquad
 b_a=\frac{K_a''(\pi/2)}{2K_a(\pi/2)}-\frac14.
\end{equation}
To justify uniformity, expand $\log\cos x=-x^2/2-x^4/12+O(x^6)$
and $K_a(\pi/2+x)$ on a fixed central interval. The remaining region
is exponentially small. At the endpoint singularity, split
$\sin^p=\sin^{p/2}\sin^{p/2}$; one factor is exponentially small,
and the other makes the inverse-linear kernel singularity uniformly
integrable for $p\ge2$.
Consequently
\[
 \frac{J_{n-k}(a)}{J_n(a)}
 =\sqrt{\frac n{n-k}}\left[1+O(k/n^2)\right]
 \quad(1\le k\le n/2).
\]
After dividing by the incoming norm and by the balanced ratio at gain
$a$, the result is at most
\[
 \frac n{\sqrt{(n-k)(n+8k)}}[1+O(k/n^2)]
 \le1-\frac{k}{4n}<1
\]
for all sufficiently large $n$. The elementary first-factor inequality
used here is $n/\sqrt{(n-k)(n+8k)}\le1-k/(2n)$.

For $k\ge n/2$, use the universal bound $|\widehat h([n])|\le\kappa_n$
for a monotone function $h:\{-1,1\}^n\to[-1,1]$.
One proof averages over coordinate permutations and writes the
resulting monotone sequence as a convex combination of Hamming-level
thresholds and constants. The threshold coefficients are signed values
$2^{1-n}\binom{n-1}{j}$, maximized at the center.
Thus the normalized coefficient is at most $\kappa_n/\sqrt{5n}$.
The balanced witness at gain one has ratio
\[
 \frac{q_n(1)}{\sqrt n}=\frac{\kappa_n}{\sqrt n}
 [F_0(1)+O(n^{-1})],\qquad
 F_0(1)\ge1-\frac4{1+e^2}>\frac12>\frac1{\sqrt5}.
\]
It is strictly better for large $n$. Coordinate sign changes affect
only the overall sign of the parity coefficient.
\end{proof}

For a fixed number of fixed odd harmonics there is a more precise loss.
If $\Delta=\sum_i(r_i^2-1)$ is fixed as $n\to\infty$, then uniformly
for $a\in[1,2]$,
\begin{equation}\label{eq:fixed-harmonic-loss}
 \frac{\left|\E\chi_{[n]}\tanh(a\sum_i r_iX_i)\right|}
 {a\sqrt{n+\Delta}}
 =\frac{q_n(a)}{a\sqrt n}
 \left[1-\frac\Delta n+O(n^{-2})\right].
\end{equation}
Indeed
\[
 P_r(y)=(-1)^{(r-1)/2}
        \frac{\sin(r\arcsin\sqrt y)}{\sqrt y}
\]
is a polynomial with $P_r(1)=1$ and $P_r'(1)=(r^2-1)/2$.
Writing $\prod_iP_{r_i}(y)=\sum_j c_jy^j$, the Fourier integral
expresses the coefficient, up to its fixed phase, as
$\sum_jc_jq_{n+2j}(a)$.
The balanced expansion gives
$q_{n+2j}/q_n=1-j/n+O(n^{-2})$ for fixed $j$.
Use $\sum_jc_j=1$ and $\sum_jjc_j=\Delta/2$, then divide by
$\sqrt{n+\Delta}$, to obtain \eqref{eq:fixed-harmonic-loss}.

\subsection{A global maximum on every fixed-dimensional defect slice}

\begin{theorem}[A fixed number of bounded real defects]
\label{thm:fixed-real-defects}
Fix an integer $k\ge1$ and $B>2$. For all sufficiently large odd $n$,
among vectors with $n-k$ coordinates equal to $a\in[1,2]$ and the
remaining $k$ coordinates in $[-B,B]$, the maximum of
\[
 \frac{|\E\chi_{[n]}(X)\tanh(u\cdot X)|}{\|u\|_2}
\]
is exactly $q_n(a_n)/(a_n\sqrt n)$. Its only maximizers have
$a=a_n$ and every exceptional coordinate equal to $a_n$ or $-a_n$.
The threshold for $n$ may depend on $k$ and $B$.
\end{theorem}
\begin{proof}
Let $r=n-k$ and $\eta=r\bmod2$. Central-binomial normalization of the
equal bulk, followed by averaging the $k$ exceptional signs, gives the
leading phase response
\[
 H(a,b)=\E_Y\chi_{[k]}(Y)
       T_a\left(\eta a+\sum_i b_iY_i\right).
\]
There is a uniform expansion through order $n^{-2}$, with every fixed
number of derivatives, on the stated compact parameter set. Here are
explicit ways to justify it. For even $r$, pair the bulk levels $\pm2j$.
The normalized expression is
\[
 \tanh c+\sum_{j\ge1}(-1)^j R_j(r)
 [\tanh(c+2aj)+\tanh(c-2aj)],\qquad
 R_j(r)=\frac{\binom r{r/2+j}}{\binom r{r/2}}.
\]
Expand $R_j(r)=1-2j^2/r+O_j(r^{-2})$ and sum the polynomial
remainders against the exponentially decreasing paired bracket.
For odd $r$, pair the levels $\pm(2j+1)$ and subtract the limiting
bracket two. The remainder and all fixed derivatives again decay
exponentially. The constant part sums exactly to the majority
central-binomial coefficient. Splitting $j$ at a small fixed multiple
of $r$, as in the proof of Theorem~\ref{thm:curvature}, makes all
remainders uniform.

The leading objective, after multiplication by $n$, is bounded by
$\sqrt{2/\pi}\,F_0(a)/a$. Theorem~\ref{thm:phase-stability} says that
equality requires every $b_i/a$ to be an odd integer. The limiting gain
has the unique maximizer $a_0\in(1,2)$ from the balanced expansion.
Its second derivative is strictly negative; this also follows by
passing the strict alternating-term estimate in the balanced-gain
proof to the limit. There are only finitely many equality patterns
in the compact box.

Near such a pattern, write $b_i=r_i a+e_i$, with odd integers $r_i$.
Up to one common phase, the leading response is
$\E T_a(a+\sum_i e_iY_i)$. Its transverse Hessian at zero is
$T_a''(a)I$, because $\E Y_iY_j=0$ for $i\ne j$.
Its gain--transverse cross derivatives vanish. Thus every leading
maximum is nondegenerate. Uniform $C^2$ convergence localizes all
finite-$n$ maximizers to these finitely many neighborhoods; strict
negative Hessians and the gradient expansion give one local maximum
in each, displaced by $O(n^{-1})$. If a pattern lies on a boundary
of the prescribed box, enlarge that local box slightly and use its
unconstrained maximum as an upper bound.

Evaluating a harmonic pattern at gain $a_n$ and using
\eqref{eq:fixed-harmonic-loss} gives the relative value
$1-\Delta/n+O(n^{-2})$. Moving to the nearby local maximum improves
this by only $O(n^{-2})$, because its gradient is $O(n^{-1})$ and its
Hessian is uniformly negative in the objective multiplied by $n$.
Every $\Delta>0$ therefore loses. For $\Delta=0$, the balanced vector
is exactly stationary and local uniqueness identifies the maximum.
\end{proof}

The fixed-dimensional theorem also gives the sharp upper bound when a
fixed number of bounded outside-coordinate weights and a bounded bias
are present, with their squared weights included in the denominator.
Apply the bias identity \eqref{eq:bias-identity} to absorb the bias
into one selected coordinate, use Theorem~\ref{thm:fixed-real-defects}
on the two resulting vectors, and use concavity of the square root.
Conditioning the finitely many outside signs and applying Jensen once
more includes their squared norm.

The estimates in this appendix are inputs to the global argument.
Appendix~\ref{app:global-single-parity} combines them with lattice
localization and residual asymptotics to handle arbitrarily many
exceptional coordinates. The code theorem then adds orientation
localization and high-order absorption.
